\chapter{Applications}

\section{Differentiation Under the Integral}

\begin{theorem}[Leibniz Rule]
    Let $f: I \times X \to \bb{R}$ where $I \subseteq \bb{R}$ is an interval. If:
    \begin{enumerate}
        \item $t \mapsto f(t, x)$ is differentiable for each $x$
        \item $|\frac{\partial f}{\partial t}(t, x)| \leq g(x)$ for some $g \in \Lp{1}(\mu)$
    \end{enumerate}
    Then $\phi(t) = \Integral_X f(t, x) d\mu(x)$ is differentiable and:
    \[
        \phi'(t) = \Integral_X \frac{\partial f}{\partial t}(t, x) d\mu(x)
    \]
\end{theorem}

\begin{proof}
    Apply DCT to the difference quotients $\frac{f(t+h, x) - f(t, x)}{h}$.
\end{proof}

\bigskip
\section{Gaussian Integral}

\begin{theorem}
    $\Integral_{\bb{R}} e^{-x^2} d\lambda(x) = \sqrt{\pi}$.
\end{theorem}

\begin{proof}
    Let $I = \Integral_\bb{R} e^{-x^2} d\lambda(x)$. Then:
    \[
        I^2 = \Integral_{\bb{R}^2} e^{-(x^2+y^2)} d(\lambda \otimes \lambda) 
        = \Integral_{[0,2\pi]} \Integral_{[0,\infty)} e^{-r^2} r\,d\lambda(r)\,d\lambda(\theta) 
        = 2\pi \cdot \frac{1}{2} = \pi
    \]
    Thus $I = \sqrt{\pi}$.
\end{proof}

